\section{Introduction}\label{sec:intro}

Non-savings, non-dividend income tax is the largest revenue devolved to the Scottish Government, and the fiscal framework makes its performance a relative question. Because the block grant deduction is indexed to the growth of the equivalent tax base in the rest of the UK (rUK), the level of Scottish revenue is immaterial and only the Scotland--rUK differential moves the budget, through the income tax net position that the Scottish Fiscal Commission (SFC) forecasts \citep{gov2023, sfc2026}. \Cref{sec:background} sets out the mechanism and its reconciliation cycle. A technology expected to shift relative earnings across regions is therefore of direct fiscal interest, and this dissertation asks whether artificial intelligence is such a technology for Scotland.

The standard input to that judgement is now an exposure score produced by a large language model, an approach \citet{eloundou2024} established and institutional forecasters have since adopted, including the \citet{obr2025} for the UK productivity outlook and \citet{henseke2025} for a UK task-based index. Recent evidence unsettles the premise that such scores are objective measurements: scoring one task set with several contemporary models moves the average exposure, and the labour-market magnitudes built on it, by multiples rather than margins \citep{yin2026}. The question this dissertation answers is what a fiscal authority can learn from LLM-based exposure measures, given that they are demonstrably model-dependent.

The question has a lineage in econometrics, and placing it there clarifies what is new. Exposure enters downstream estimates as a regressor measured with error, the classical setting in which noise attenuates and systematic bias contaminates \citep{bound1991}. % TODO(bib): add Bound & Krueger (1991), or Bound, Brown & Mathiowetz (2001), to references.bib under the key cited here.
The regional index is a shift-share object, so the algebra of employment-share weights governs how any scoring error aggregates. And the central result works by differencing: the Scotland--rUK differential differences out the model fixed effect akin to how difference-in-differences differences out a unit fixed effect.

Three contributions follow, one for each of the questions posed at the end of \cref{sec:background}. The first is analytical. Proposition 1 shows that the regional differential in an employment-weighted exposure index removes the common component of any model-specific bias, retaining only the part co-varying with the cross-region employment-share gap, and Corollary 1 bounds that part by a Cauchy--Schwarz inequality the data can be checked against. Corollary 2 covers the case the data turn out to present, in which a model disagrees about the scale of the index while agreeing about its origin, and there a standardised or rank-based differential removes the bias in full. I take the result to the data by re-scoring the full task set on two further language models. The second contribution is substantive: a compositional answer for a devolved fiscal authority, locating Scotland fifth of the twelve UK regions and nations, attributing about two-fifths of the headline gap to London, and identifying a tilt towards the substitution channel in the composition of Scottish exposure. The third is a propagation template that carries scoring uncertainty into the objects a fiscal body forecasts, from the aggregate index through a TFP functional and a micro wage estimate to a stylised net-position sensitivity. The second and third are contributions rather than applications because the compositional weighting works at two channels: the same structure that zeroes the within-occupation term of the shift-share decomposition also demeans the model bias. 


The findings preview as follows. Scotland's employment-weighted continuous exposure sits 0.86 index points below the rUK on the 0--100 scale, with a composed 95\% interval of $[-1.07, -0.69]$ index points, most of whose width comes from survey sampling. Scale-free, that is about two percentile points of the UK occupational exposure distribution. The sign survives two full re-scores of all 17,945 tasks, the raw gap moving to $-0.53$ under Claude Haiku and $-0.43$ under GPT-4o as absolute levels shift by up to 9.1 index points; most of that movement is the scale compression Corollary 2 describes, and the residual is 23--24\% of the Cauchy--Schwarz worst case. The classified exposed share, the object closest to the policy literature's headline numbers, is censored at 91\% of employment and its sign is unstable to its own thresholds, so it cannot carry the regional question. Passed through the Acemoglu calibration, the differential implies a ten-year TFP differential of under three basis points, which \cref{sec:fiscal} translates into a net-position sensitivity of about $£1.5$ million a year.

The rest of the paper is organised as follows. \Cref{sec:background} sets out the fiscal framework, the task-based mechanism, and the measurement literature, ultimately proposing the three research questions I address. \Cref{sec:data} describes the O*NET task data, the crosswalk to UK SOC 2020, and the pooled Annual Population Survey employment and earnings data. \Cref{sec:method} constructs the exposure scores, states the region-invariance assumption, proves the cancellation results, and sets out the calibration and propagation design. \Cref{sec:results} presents the differential, its anatomy, its robustness, and its translation into productivity, wages, and the net position. \Cref{sec:conclusion} concludes.